\documentclass[11pt,a4paper]{article}
\usepackage[margin=18mm,headheight=14pt]{geometry}
\usepackage[T1]{fontenc}
\usepackage{lmodern}
\usepackage{amsmath,amssymb}
\usepackage{booktabs}
\usepackage{graphicx}
\usepackage{xcolor}
\usepackage{hyperref}
\usepackage{fancyhdr}
\usepackage{enumitem}
\usepackage{microtype}
\usepackage{siunitx}
\hypersetup{colorlinks=true,linkcolor=blue!55!black,urlcolor=blue!55!black}
\definecolor{noteBlue}{HTML}{2166AC}
\definecolor{noteRed}{HTML}{B2182B}
\definecolor{noteGray}{HTML}{4D4D4D}
\newcommand{\Usym}{U_{\rm sym}}
\newcommand{\Ekin}{E_{\rm kin}}
\newcommand{\HGND}{HGND}
\newcommand{\Rhard}{R=N(\Ekin>2\,\mathrm{GeV})/N(\Ekin<1\,\mathrm{GeV})}
\newcommand{\notegraphic}[2][]{%
  \IfFileExists{#2}{\includegraphics[#1]{#2}}{%
    \fbox{\parbox[c][42mm][c]{0.88\linewidth}{\centering
      Figure file: \texttt{\detokenize{#2}}}}}}
\pagestyle{fancy}
\fancyhf{}
\lhead{BM@N HGND internal analysis note}
\rhead{29 September 2026}
\cfoot{\thepage}
\setlength{\parindent}{0pt}
\setlength{\parskip}{5pt}

\begin{document}

\begin{center}
{\Large\bfseries Feasibility of neutron--proton observables and spectral hardness\\
with the BM@N Highly-Granular Neutron Detector}\par
\vspace{4mm}
{\large Full-production truth study and reconstruction update}\par
\vspace{3mm}
{\color{noteRed}\bfseries Internal analysis note -- not for external distribution}\par
\end{center}

\begin{abstract}
Three SMASH Xe+Cs samples with symmetry-potential strengths
$\Usym=0$, 18 and 90~MeV provide 2,345,993 events in 568 independent
generation jobs.  This note begins from the observable question.  The
integrated neutron-to-proton yield ratio has differential kinematic structure,
but it is effectively insensitive to $\Usym$ in the exact HGND front-face
acceptance.  Spectral hardness is different: the neutron hardness rises by
$4.52\%$ and the proton hardness falls by $5.39\%$ between 0 and 90~MeV,
giving a $10.48\%$ truth-level double-ratio response.  These conclusions use
impact-parameter reweighting and a bootstrap over production jobs.

The reconstructed observable must be defined inside the detector-efficiency
plateau.  With one frozen GNN and the same score threshold for all three SMASH
samples, neither the raw neutron-candidate yield nor the detector-safe neutron
hardness separates the 0 and 90~MeV samples significantly.  A hybrid
reconstructed-neutron/truth-proton hardness reaches only $1.7\sigma$, compared
with $13.8\sigma$ for the ideal truth double ratio, and is a feasibility test
rather than a measurement.  An event-level study quantifies the efficiency for
the first, second and subsequent true-neutron clusters rather than quoting only
an inclusive average.
\end{abstract}

\section{Decision summary}

\begin{center}
\begin{tabular}{p{0.29\linewidth}p{0.25\linewidth}p{0.37\linewidth}}
\toprule
Question & Result & Analysis decision \\
\midrule
Integrated $n/p$ in HGND & $+0.06\%\pm0.28\%$ for 90/0 & Do not use as a
symmetry-potential estimator. \\
Differential $n/p(\Ekin,y,p_T)$ & Strong shape dependence & Retain as a
transport-model and phase-space diagnostic. \\
Neutron hardness in HGND & $+4.52\%\pm0.45\%$ & Primary detector-level
physics target. \\
Proton hardness in same window & $-5.39\%\pm0.51\%$ & Requires the charged
tracking arm, not HGND alone. \\
Hardness double ratio & $+10.48\%\pm0.76\%$ & Preferred truth-level
isovector diagnostic. \\
Conventional reconstructed $R_n$ & Denominator below 1~GeV has $\sim2\%$
efficiency & Redefine the reconstructed window inside the efficiency plateau. \\
Detector-safe reconstructed $\widehat R_n$ & $-0.39\%\pm2.86\%$ for 90/0 &
Current reconstruction is statistically compatible with no response. \\
Modified 1--2~GeV upper-tail hardness & $+1.49\%\pm1.39\%$ for 90/0 &
Best ordered in-band ratio, but exploratory and only $1.07\sigma$. \\
\bottomrule
\end{tabular}
\end{center}

\section{Neutron-to-proton ratio feasibility}

\subsection{Global phase space}

The global primary-nucleon ratio reproduces the qualitative structure of the
``n/p--SMASH check'' reference: $n/p$ separates most strongly at high kinetic
energy and its shape depends on rapidity.  Figure~\ref{fig:globalnp} uses the
complete available production rather than a file subset.  The lower panels
show the ratio to the $\Usym=0$ sample so that the parameter response is not
confused with the large common spectral shape.

\begin{figure}[htbp]
\centering
\notegraphic[width=\linewidth]{../results/update_note/fig01_np_global_ekin.pdf}
\caption{Primary-nucleon $n/p$ versus kinetic energy at midrapidity (left) and
in full phase space (right).  All available production files are included.
Statistical errors are evaluated with event clustering in the source reduction.}
\label{fig:globalnp}
\end{figure}

The earlier right-hand panel used the disjoint selection
$|y-y_{\rm cm}|>0.5$.  Its two kinematic branches exchange dominance near
1.1~GeV, producing an abrupt, visually misleading structure.  The revised
panel shows inclusive rapidity; all valid points are connected consistently in
the spectrum and ratio panels.  The global differential separation is real but is not by itself an HGND
measurement.  HGND covers a rectangular front face at negative $x$, only
0.0114~sr, and the SMASH reaction plane is fixed.  A polar-angle-only selection
therefore cannot replace the exact $(x,y)$ acceptance.

\subsection{Exact HGND front-face acceptance}

For a sample $s$, event weights $w_e$, a bin $B=[x_a,x_b)$ of width
$\Delta x_B=x_b-x_a$, and particle species $a\in\{n,p\}$, the differential
per-event yield used in this note is
\[
 Y_a^{(s)}(B)\equiv\left.\frac{dN_a}{dx}\right|_B
 =\frac{\sum_e w_e N_{a,e}^{(s)}(B)}
        {\Delta x_B\sum_e w_e}.
\]
The differential neutron-to-proton yield ratio is
\[
 \rho_{n/p}^{(s)}(B)\equiv\frac{Y_n^{(s)}(B)}{Y_p^{(s)}(B)}
 =\frac{\sum_e w_eN_{n,e}^{(s)}(B)}
        {\sum_e w_eN_{p,e}^{(s)}(B)},
\]
where the second equality is valid only when neutron and proton yields use the
same events, weights, bin and acceptance.  An integrated yield is obtained by
replacing $B$ by the declared phase-space region and omitting $\Delta x_B$.
For any observable $O$, the quoted Spot response is
$\Delta_O^{90/0}=O^{(90)}/O^{(0)}-1$.

For selected reconstructed neutron candidates the corresponding raw yield is
\[
 \widehat Y_n^{(s)}(B)=
 \frac{\sum_c w_{e(c)}\,\mathbf{1}(q_c>q_0)\,
       \mathbf{1}(x_c^{\rm reco}\in B)}
      {\Delta x_B\sum_e w_e}.
\]
It includes fake candidates unless a response model is applied.  The hybrid
$\widehat\rho_{n/p}=\widehat Y_n/Y_p^{\rm MC}$ therefore tests feasibility
with MC-truth protons; it is not a data-level $n/p$ measurement.

Figure~\ref{fig:accnp} propagates both neutron and proton primaries from their
production vertices to the detector front face.  The uncertainty is obtained by
resampling whole generation files; treating all particles as independent would
substantially underestimate it.  The energy-differential ratio separates above
roughly 2~GeV, but its integral over the acceptance does not order the three
$\Usym$ values.  The integrated 90/0 change is only $+0.06\%$ ($0.2\sigma$).
The apparent 2--3~GeV gap in the previous rendering was not physical.  It came
from assigning logarithmic native bins to coarse bins by their centres.  The
revised plot merges complete native intervals using their exact boundaries,
conserves all entries and connects all populated points.

\begin{figure}[htbp]
\centering
\notegraphic[width=0.78\linewidth]{../results/update_note/fig02_np_hgnd_acceptance_ekin_v2.pdf}
\caption{Differential primary-nucleon $n/p$ in the exact HGND front-face
acceptance.  The lower panel is the 90/0 double ratio.  Error bars are from a
400-replica generation-job bootstrap.}
\label{fig:accnp}
\end{figure}

The conclusion is deliberately narrow: integrated $n/p$ is not a useful
$\Usym$ observable for HGND.  Differential $n/p$ remains valuable for generator
validation and for a combined neutron--charged-particle analysis with matched
phase space.

The angularly differential neutron yield and $n/p$ response provide a useful
cross-check because they retain information that cancels in the integral.  In
the HGND angular band, however, the truth neutron-yield response is only
$+0.181\%\pm0.153\%$ ($1.18\sigma$).  The much larger changes below about
$6^\circ$ are outside the HGND band and cannot be claimed as HGND sensitivity.

\begin{figure}[htbp]
\centering
\notegraphic[width=0.92\linewidth]{../results/update_note/fig07_differential_neutron_yield_np_theta.pdf}
\caption{Relative per-event neutron yield (top) and $n/p$ response (bottom)
versus polar angle.  The HGND angular band is shaded; errors are
event-clustered.}
\label{fig:differential-yield}
\end{figure}

\section{Spectral hardness is the recommended observable}

For the truth study we retain
\[
  \Rhard.
\]
It removes the overall yield normalisation and tests the energy redistribution
that the symmetry potential produces.  The 18~MeV point lies between the two
outer samples for $R_n$, $R_p$ and $R_n/R_p$.

\begin{figure}[htbp]
\centering
\notegraphic[width=0.78\linewidth]{../results/update_note/fig03_neutron_hardness_global_vs_hgnd.pdf}
\caption{Neutron spectral hardness globally and in the HGND front-face
acceptance.  The response is localised in the forward detector window; the
global spectrum changes little.}
\label{fig:hardness}
\end{figure}

\begin{figure}[htbp]
\centering
\notegraphic[width=\linewidth]{../results/update_note/fig04_observable_feasibility_summary.pdf}
\caption{Observable comparison.  Yield $n/p$ is small or null in the HGND
acceptance, while the spectral-hardness observables show a monotonic isovector
response.}
\label{fig:summary}
\end{figure}

At full statistics the exact-acceptance responses are
\begin{align*}
  R_n(90)/R_n(0)-1 &= +4.52\%\pm0.45\%,\\
  R_p(90)/R_p(0)-1 &= -5.39\%\pm0.51\%,\\
  \frac{(R_n/R_p)_{90}}{(R_n/R_p)_0}-1 &= +10.48\%\pm0.76\%.
\end{align*}
The larger angular band gives $+10.73\%$ for the last quantity.  The agreement
shows that the response is not created by the fiducial definition; the exact
acceptance simply has 7.9 times fewer particles and hence larger uncertainty.

The double-ratio response is not caused solely by protons.  In logarithmic
response, the neutron hardening contributes
$\ln(1.0452)=0.0442$ and the proton softening contributes
$-\ln(0.9461)=0.0554$: approximately 44\% and 56\% of the total, respectively.
The neutron spectra can look nearly coincident on a logarithmic yield plot
while their high/low integrals still differ at more than $10\sigma$ because the
full production sample makes a few-percent shape change statistically precise.

\subsection{Separation power and hybrid feasibility estimates}

Table~\ref{tab:separation} and Fig.~\ref{fig:separation} compare the 90/0
response using $|\Delta|/\sigma_\Delta$ as a common separation metric.  Every
reconstructed entry is derived from the same frozen GNN, common feature scaler
and common selection $q>0.45$ on 409,600 represented events per Spot sample.
The raw candidate quantities include fake clusters; the matched yield is shown
only as a closure diagnostic.

Three additional reconstructed shape observables test whether the populated
1--2~GeV interval can provide a more precise hardness measurement:
\begin{align*}
 R_{n,{\rm eq}}^{1-2}
 &=\frac{N(1.5\leq E_{\rm reco}<2.0)}
          {N(1.0\leq E_{\rm reco}<1.5)},\\
 R_{n,{\rm tail}}^{1-2}
 &=\frac{N(1.8\leq E_{\rm reco}<2.0)}
          {N(1.0\leq E_{\rm reco}<1.8)},\\
 \langle E\rangle_{1-2}
 &=\frac{\sum_c E_c^{\rm reco}\,
          \mathbf{1}(1.0\leq E_c^{\rm reco}<2.0)}
         {\sum_c\mathbf{1}(1.0\leq E_c^{\rm reco}<2.0)}.
\end{align*}
All sums use selected candidates with $q>0.45$.  These are raw,
detector-folded quantities rather than efficiency-corrected spectra.
For each shape definition
$x\in\{{\rm eq},{\rm tail},\langle E\rangle\}$, the corresponding hybrid is
\[
 H_x=\frac{O_{n,x}^{\rm reco}}{O_{p,x}^{\rm MC}},\qquad
 \Delta_{H_x}^{90/0}=\frac{(H_x)_{90}}{(H_x)_0}-1.
\]
The proton quantities use primary MC protons intersecting the exact HGND front
face.  Their uncertainties come from a 400-replica production-file bootstrap.

\begin{table}[htbp]
\centering
\small
\begin{tabular}{lrrl}
\toprule
Observable & Response [\%] & Separation & Interpretation \\
\midrule
Truth differential $n$ yield & $+0.181\pm0.153$ & $1.18\sigma$ & weak in HGND band \\
Truth yield $n/p$ & $+0.062\pm0.283$ & $0.22\sigma$ & not sensitive \\
Truth $R_n$ & $+4.519\pm0.445$ & $10.15\sigma$ & neutron contribution \\
Truth $R_p$ & $-5.391\pm0.515$ & $10.48\sigma$ & proton contribution \\
Truth $R_n/R_p$ & $+10.475\pm0.758$ & $13.82\sigma$ & ideal combined observable \\
Reco candidate $\widehat Y_n$ & $+0.418\pm0.467$ & $0.89\sigma$ & includes fakes \\
Reco matched $\widehat Y_n^{\rm match}$ & $+0.905\pm0.511$ & $1.77\sigma$ & closure diagnostic \\
Reco $d\widehat N_n/dE$ & $+1.727\pm2.909$ & $0.59\sigma$ & 2.6--3.0~GeV \\
Reco safe $\widehat R_n$ & $-0.391\pm2.863$ & $0.14\sigma$ & detector-safe windows \\
Reco $R_{n,{\rm eq}}^{1-2}$ & $+1.088\pm1.284$ & $0.85\sigma$ & not ordered \\
Reco $R_{n,{\rm tail}}^{1-2}$ & $+1.494\pm1.395$ & $1.07\sigma$ & ordered, exploratory \\
Reco $\langle E\rangle_{1-2}$ & $+0.043\pm0.106$ & $0.41\sigma$ & not ordered \\
Hybrid $\widehat Y_n/Y_p^{\rm MC}$ & $+0.149\pm0.492$ & $0.30\sigma$ & provisional \\
Hybrid $\widehat R_n/R_p^{\rm MC}$ & $+5.285\pm3.080$ & $1.72\sigma$ & provisional \\
Hybrid $R_{n,{\rm eq}}^{1-2}/R_{p,{\rm eq}}^{1-2,{\rm MC}}$ & $+1.916\pm1.480$ & $1.29\sigma$ & equal-bin shape \\
Hybrid $R_{n,{\rm tail}}^{1-2}/R_{p,{\rm tail}}^{1-2,{\rm MC}}$ & $+2.799\pm1.677$ & $1.67\sigma$ & upper-tail shape \\
Hybrid $\langle E_n\rangle_{1-2}/\langle E_p\rangle_{1-2}^{\rm MC}$ & $+0.144\pm0.124$ & $1.17\sigma$ & mean-energy shape \\
\bottomrule
\end{tabular}
\caption{Separation power of truth, reconstructed-neutron and hybrid
reconstructed-neutron/truth-proton observables.  Hybrid uncertainties use an
independent-error approximation and therefore are feasibility estimates only.}
\label{tab:separation}
\end{table}

\begin{figure}[htbp]
\centering
\notegraphic[width=0.94\linewidth]{../results/update_note/fig06_observable_separation_power_v5.pdf}
\caption{Relative 90/0 response and separation power.  Reconstruction strongly
reduces the neutron-only sensitivity.  All reconstructed-neutron and hybrid
points use one frozen neutron model and one score threshold.}
\label{fig:separation}
\end{figure}

The hybrid results establish feasibility, not final sensitivity.  The current
neutron and proton products do not share an event-by-event sample, and the
quoted uncertainty assumes independent inputs.  A publishable result requires
reconstructed charged tracks in the same front-face-equivalent phase space and
a common event/job bootstrap, including neutron--proton covariance.  No current
reconstructed quantity reaches $2\sigma$; the apparent truth-level advantage
of hardness is not yet retained by reconstruction.

For the new proton components, the 90/0 responses are
$-0.812\%\pm0.698\%$ for the equal-bin hardness,
$-1.269\%\pm0.869\%$ for the upper-tail hardness, and
$-0.101\%\pm0.064\%$ for the in-band mean energy.  Combining the opposite-sign
neutron and proton changes raises the upper-tail hybrid to $1.67\sigma$, close
to but not better than the $1.72\sigma$ conventional hybrid.  The neutron and
proton samples are not event matched, so all hybrid errors use independent
propagation and omit neutron--proton covariance.

The full-production primary CSVs used for these exact proton windows contain
the negative impact-parameter sentinel $B=-1$.  Consequently the new proton
components are unweighted full-production results; no centrality reweighting
is claimed for them.  Their file/event scope is nevertheless complete:
195/200/173 files and 804,369/824,338/717,286 events for 0/18/90~MeV.

The 1--2~GeV ratios nevertheless improve the reconstructed statistical
precision.  The equal-bin definition has a non-monotonic central sequence
$1.7279/1.7100/1.7467$ for 0/18/90~MeV.  The upper-tail definition is ordered,
$0.3694/0.3707/0.3749$, but was selected after examining several boundaries;
its nominal $1.07\sigma$ separation is therefore descriptive rather than a
blind significance.  It should be frozen on development data and evaluated on
an independent sample.  A forward-folded binned shape fit over 1--2~GeV
remains preferable because it uses the full spectrum and explicitly models
efficiency and energy migration.

\subsection{SMASH reconstruction performance}

The same model was evaluated on 1,831,257 SMASH cluster rows across the three
Spot samples.  Its cluster-classification ROC AUC is 0.956.  At the fixed
$q>0.45$ working point, count purity is 84.7\% and count efficiency is 64.9\%
(energy-weighted purity and efficiency are 86.4\% and 72.0\%).  Energy
resolution improves from about 21\% below 1~GeV to 9\% near 3.75~GeV, but the
median energy response changes sign near 2~GeV and reaches a $-20\%$ bias at
2.8~GeV.  This compression explains why a truth threshold cannot simply be
copied into reconstructed energy.

\begin{figure}[htbp]
\centering
\notegraphic[width=\linewidth]{../results/update_note/fig08_smash_reconstruction_performance.pdf}
\caption{SMASH reconstruction performance for one model: purity--efficiency,
median energy bias and robust energy resolution.  The star marks the common
score working point used by every reconstructed observable in
Fig.~\ref{fig:separation}.}
\label{fig:recoperf}
\end{figure}

\subsection{Centrality}

The production has job-level impact-parameter block structure.  The three full
samples nevertheless have compatible means,
$\Delta\langle b\rangle_{90-0}=+0.0034\pm0.0147$~fm.  Event-level reweighting
reduces the residual spread in the three means to $9.5\times10^{-5}$~fm and
costs 0.1\% in effective statistics.  Figure~\ref{fig:cent} demonstrates that
the isovector response persists across centrality; the yield ratio does not
provide a stable alternative.

\begin{figure}[htbp]
\centering
\notegraphic[width=0.82\linewidth]{../results/update_note/fig05_hgnd_response_vs_centrality.pdf}
\caption{Truth-level 90/0 response in ten impact-parameter classes after
reweighting.  Every uncertainty is evaluated by resampling generation jobs and
rebuilding the weights in each replica.}
\label{fig:cent}
\end{figure}

\section{Reconstruction efficiency and neutron multiplicity}

An inclusive cluster efficiency is insufficient when the physics sample
contains two or more neutrons in the same event.  The relevant questions are
conditional:
\begin{align*}
 \varepsilon_k &\equiv P(\text{$k$th-hardest true neutron selected}
                  \mid N_n^{\rm true}\geq k),\\
 C_k &\equiv P(N_n^{\rm selected,true}\geq k
                  \mid N_n^{\rm true}\geq k).
\end{align*}
Truth here means a truth-matched reconstructable neutron cluster
($\texttt{cl\_label}=1$), and true neutrons are ranked by decreasing true
kinetic energy.  A single classifier working point is applied to all three
samples.  The resulting rank efficiencies and the row-normalised multiplicity
response are shown in Figs.~\ref{fig:rankeff} and~\ref{fig:multresponse}.

\begin{figure}[htbp]
\centering
\notegraphic[width=\linewidth]{../results/update_note/conditional_neutron_efficiency_v2.pdf}
\caption{Conditional efficiency of the first through fourth true-neutron
clusters and efficiency versus true event multiplicity.  Wilson intervals are
statistical only, and plotted points require at least 20 eligible events.  The
dashed curves require all neutrons through rank $k$ to be retained and expose
compounding losses in multi-neutron events.}
\label{fig:rankeff}
\end{figure}

\begin{figure}[htbp]
\centering
\notegraphic[width=\linewidth]{../results/update_note/multiplicity_response.pdf}
\caption{Row-normalised response between reconstructable truth multiplicity and
selected truth-matched multiplicity.  The last bin includes overflow.  Fake
selected clusters are excluded here and are tracked separately through purity.}
\label{fig:multresponse}
\end{figure}

At the common score threshold 0.45, the first-neutron rank efficiency is
66.6--66.8\% across the three Spot samples.  Conditional on a second true
neutron it is 46.4--47.7\%, and for a third it is 36.5--41.1\%.  The probability
of retaining all truth-matched neutrons is only 31.9--33.7\% for true
multiplicity two and 14.1--14.7\% for multiplicity three.  Fourth-neutron rank
estimates are sample limited (24--35 eligible events) and are not yet suitable
for architecture ranking.

Within these uncertainties the three datasets are consistent for ranks 1--3;
this is a response-closure check, not proof of identical efficiencies.  The
former 90~MeV point at rank 5 (the horizontal axis is neutron rank, not energy)
contained exactly one eligible event.  It is retained in the CSV/JSON audit
products but omitted from the physics figure by the predeclared
$N_{\rm event}\geq20$ display requirement.

This decomposition should replace a single ``neutron efficiency'' number in
algorithm comparisons.  The first-neutron efficiency mainly measures ordinary
acceptance and classification.  The fall of $\varepsilon_k$ or $C_k$ with rank
isolates overlap, cluster merging and low-energy losses.  The same curves must
be shown for the baseline and every challenger architecture at a common purity
or with their full purity--efficiency envelopes.

These conditional probabilities are detector-response inputs, not candidate
Spot observables.  Their purpose is to correct the neutron yield and energy
spectrum.  In matrix form, the selected multiplicity/energy distribution is
$m_i=\sum_j A_{ij}t_j+b_i$, where $A_{ij}$ includes the rank-dependent losses
and migration and $b_i$ is the fake-cluster contribution.  A forward-folded fit
of $t_j$, or a regularised inversion used only after closure tests, can recover
neutrons lost preferentially in multiplicity-two and multiplicity-three
events.  This matters for Spot sensitivity even when the efficiency curves are
identical across samples: if Spot changes the joint energy--multiplicity
composition, applying only one inclusive efficiency biases both the corrected
yield and the high/low hardness ratio.  The response therefore has to be binned
at least in true/reconstructed multiplicity and neutron energy, with its
uncertainty propagated through the 90/0 comparison.

\section{Energy response and observable definition}

The efficiency for neutrons reaching the front face varies from 0.02 below
0.5~GeV and begins its rapid rise only around 0.5--0.7~GeV, reaching
approximately 0.80 around 2--3~GeV.  The measured threshold ratio is therefore
not the truth ratio but
\[
 R_{\rm obs}=\frac{\int_H \epsilon(E)f(E)\,dE}
                   {\int_L \epsilon(E)f(E)\,dE}.
\]
For $L:\Ekin<1$~GeV, the denominator is dominated by the accepted
0.7--1.0~GeV tail rather than by the full truth population.  Direct inverse
efficiency weights reach roughly 50 below 0.5~GeV, creating extreme variance
and sensitivity to small modelling errors.  Energy migration through the
1~GeV boundary adds a second bias, and cancellation between Spot samples is
incomplete because their spectra are not identical where $\epsilon(E)$ is
steep.  Thus the turn-on can attenuate, amplify or even reverse the measured
90/0 response; it is not merely an overall loss of events.

Consequently the denominator of the truth definition $\Rhard$ is not robustly
recoverable by bin-by-bin unfolding.  The reconstructed hardness observable
must either use two energy regions inside the validated efficiency plateau
(the current held-out study uses 1.8--2.4 and $>2.6$~GeV) or be extracted from
a forward-folded spectral fit with efficiency, energy-scale and migration
nuisance parameters.  Thresholds are selected on training+validation only and
frozen before the test sample is opened.

The common feature normalisation is mandatory.  Fitting a separate scaler in
each sample changed the scale of the time-of-flight energy feature by 1.9\% and
removed part of the physical difference before inference.  Reusing one scaler
reduced the differential high-energy response from
$-30.7\pm8.6$ to $-11.4\pm9.8$~MeV and restored the sign and ordering of the
reconstructed central values without changing the network weights.

\section{Algorithm comparison and release criteria}

The paper-facing analysis should describe a frozen baseline and one or more
challenger architectures, not the chronological history of training.  Every
model must use the same job-level train/calibration/test split, shared scaler,
truth definition, front-face selection and pre-declared working points.  The
minimum comparison is:
\begin{enumerate}[leftmargin=*]
\item cluster ROC and purity--efficiency curve;
\item energy linearity and resolution versus true energy;
\item $\varepsilon_k$ and $C_k$ through at least the fourth neutron;
\item the complete multiplicity response matrix;
\item reconstructed hardness on a blind, adequately powered test set;
\item stability under seed, score threshold and energy-window variation.
\end{enumerate}

A challenger passes the physics gate only if reconstructed and truth 90/0
effects have the same sign, their paired difference is below $2\sigma$ and
below 25\% of the truth effect, the differential response is consistent with
zero, and the 0/18/90 ordering survives.  If no architecture passes, the
measurement should use a forward-folded spectral fit with response nuisance
parameters instead of an unfolded threshold ratio.

\section{Current internal status and next actions}

The resumable end-to-end CUDA training has a validated checkpoint after ten
epochs; the best validation loss is 0.7349 at epoch 9.  The resume-state defect
that stopped the continuation was corrected and covered by a regression test,
and a test-partition resume completed successfully before the rocky job was
relaunched.  These operational details belong in this update note and are not
paper content.

The remaining ordered work is:
\begin{enumerate}[leftmargin=*]
\item train the challenger to the frozen stopping condition;
\item evaluate all three samples and regenerate the multiplicity and response
figures from that checkpoint;
\item repeat the same plots for the baseline and at least one architecture
challenger;
\item freeze the blind test before any reconstructed sensitivity statement;
\item move only the frozen baseline--challenger comparison into the paper.
\end{enumerate}

\section*{Reproducibility record}

The truth reductions contain 195/200/173 generation files and
804,369/824,338/717,286 events for $\Usym=0/18/90$~MeV.  Plot inputs and
job-bootstrap products are stored under
\path{results/np_ratio_smash_check_full}, \path{results/acceptance_full}, and
\path{results/update_note}.  The reconstructed comparison contains 409,600
represented events per Spot sample and is recorded in
\path{results/update_note/smash_reconstructed_observables.json}.  The figure builder is
\path{scripts/make_update_note_figures.py}; the multiplicity analysis is
\path{scripts/multiplicity_efficiency.py}.

\end{document}
